\subsection{LIMIT POINT AND CLUSTER POINT OF A FUNCTION}

DEFINITION 3. Let \(f\) be a mapping of a set \(\mathbf{X}\) into a topological space \(\mathbf{Y}\) , and let \(\mathfrak{F}\) be a filter on \(\mathbf{X}\) . A point \(y \in \mathbf{Y}\) is said to be a limit point (or simply a limit) (resp. cluster point) of \(f\) with respect to the filter \(\mathfrak{F}\) if \(y\) is a limit point (resp. cluster point) of the filter base \(f(\mathfrak{F})\) .

The relation ``y is a limit of f with respect to the filter \(\mathfrak{F}\) '' is written \(\lim_{\mathfrak{F}} f = y\) , or \(\lim_{x, \mathfrak{F}} f(x) = y\) , or \(\lim_{x} f(x) = y\) if there is no risk of confusion.

From Definition 3 and Propositions 1 (no. 1) and 3 (no. 2) we deduce the following criteria :

PROPOSITION 7. A point \(y \in \mathbf{Y}\) is a limit of \(f\) with respect to the filter \(\mathfrak{F}\) if and only if, for each neighbourhood \(\mathbf{V}\) of \(y\) in \(\mathbf{Y}\) , there is a set \(\mathbf{M} \in \mathfrak{F}\) such that \(f(\mathbf{M}) \subset \mathbf{V}\) (i.e.~\(\overline{f}^{1}(\mathbf{V}) \in \mathfrak{F}\) for each neighbourhood \(\mathbf{V}\) of \(y\) ).

A point \(y \in Y\) is a cluster point of f with respect to F if and only if for each neighbourhood V of y and each \(M \in F\) there is a point \(x \in M\) such that \(f(x) \in V\) .

Examples. 1) A sequence of points \((x_{n})_{n\in \mathbb{N}}\) of a topological space is a mapping \(n\to x_n\) of \(\mathbf{N}\) into \(\mathbf{X}\) . In analysis one frequently uses the notions of limit point and cluster point of such a mapping with respect to the Fréchet filter (§ 6, no. 1) on \(\mathbf{N}\) ; if \(y\) is a limit of \(n\to x_{n}\) with respect to the Fréchet filter, \(y\) is said to be a limit of the sequence \((x_{n})\) as \(n\) tends to infinity, and we write \(\lim_{n\to \infty}x_n = y\) . A cluster point of the mapping

\(n \to x_{n}\) with respect to the Fréchet filter is called a cluster point of the sequence \((x_{n})\) .

Thus a point \(y \in \mathbf{X}\) is a limit (resp. cluster point) of a sequence \((x_n)\) of points of \(\mathbf{X}\) if it is a limit point (resp. cluster point) of the elementary filter associated with \((x_n)\) ( \(\S 6\) , no. 8).

The point y is a limit of a sequence \((x_{n})\) in X if and only if, for every neighbourhood V of y in X, all but a finite number of the terms of the sequence \((x_{n})\) are in V, i.e.~there is an integer \(n_{0}\) such that \(x_{n} \in V\) for all \(n \geqslant n_{0}\) . Likewise y is a cluster point of the sequence \((x_{n})\) if and only if, for every neighbourhood V of y in X and every integer \(n_{0}\) , there is an integer \(n \geqslant n_{0}\) such that \(x_{n} \in V\) .

\begin{enumerate}
\def\labelenumi{\arabic{enumi})}
\setcounter{enumi}{1}
\tightlist
\item
  More generally, let \(f\) be a mapping of a directed set \(A\) into a topological space \(X\) . If \(x \in X\) is a limit (resp. cluster point) of \(f\) with respect to the section filter of \(A\) then \(x\) is said to be a limit (resp. cluster point) of \(f\) with respect to the directed set \(A\) , and we write \(x = \lim f(z)\) .
\end{enumerate}

If y is a limit (resp. cluster point) of a mapping \(f: X \rightarrow Y\) with respect to a filter F on X, then y remains a limit (resp. cluster point) of f with respect to F if we replace the topology of Y by a coarser topology, or if we replace the filter F by a finer (resp. coarser) filter.

PROPOSITION 8. Let \(f\) be a mapping of a set \(X\) into a topological space \(Y\) ; then \(y \in Y\) is a cluster point of \(f\) with respect to \(\mathfrak{F}\) if and only if there is a filter \(\mathfrak{G}\) on \(X\) which is finer than \(\mathfrak{F}\) and such that \(y\) is a limit of \(f\) with respect to \(\mathfrak{G}\) .

For if y is a cluster point of f with respect to F, and if B is the neighbourhood filter of y, then \(\overline{f}^{1}(\mathfrak{B})\) is a filter base on X since every set of \(\overline{f}^{1}(\mathfrak{B})\) meets every set of F (§ 6, no. 6). This remark shows also that there is a filter G on X which is finer than both F and the filter with base \(\overline{f}^{1}(\mathfrak{B})\) (§ 6, no. 2, Proposition 1, Corollary 2), hence that y is a limit point of f with respect to G.

Notice finally that if f is a mapping of a set X into a topological space Y, the set of cluster points of f with respect to a filter F on X is closed in Y (no. 2, Proposition 5) and possibly empty.

Remark. If \(y \in Y\) is a limit (resp. cluster point) of a mapping \(f: X \to Y\) with respect to a filter F on X, then y is also a limit (resp. cluster point) of every function \(g: X \to Y\) which has the same germ as f with respect to F (§ 6, no. 9); y is said to be a limit (resp. cluster point) of the germ \(\tilde{f}\) of f with respect to F.

